% ECONOMICS_PROBLEM: {"id": "F37-573", "title": "Exchange-rate forecasting", "jel_code": "F37", "source_id": "OP-0330"}
% !TeX program = xelatex
\documentclass[11pt,a4paper]{article}
\usepackage[margin=23mm]{geometry}
\usepackage{fontspec}
\setmainfont{Latin Modern Roman}
\usepackage{amsmath,amssymb,mathtools}
\usepackage{xcolor,enumitem,booktabs,longtable,array,xurl,hyperref}
\definecolor{muted}{HTML}{53636C}
\hypersetup{colorlinks=true,urlcolor=blue,linkcolor=blue}
\setlength{\parindent}{0pt}
\setlength{\parskip}{5pt}
\newcommand{\R}{\mathbb R}
\newcommand{\E}{\mathbb E}
\newcommand{\N}{\mathbb N}
\newcommand{\ind}{\mathbf 1}
\newcommand{\Var}{\operatorname{Var}}
\newcommand{\Cov}{\operatorname{Cov}}
\newcommand{\supp}{\operatorname{supp}}
\DeclareMathOperator*{\argmin}{arg\,min}
\DeclareMathOperator*{\argmax}{arg\,max}
\newcommand{\entrymeta}[3]{{\sffamily\small\color{muted}\textbf{\texttt{#1}}\quad|\quad #2\par #3\par}}
\newcommand{\Problemref}[1]{\texttt{#1}}
\newcommand{\JELref}[2]{\texttt{#2} (JEL \texttt{#1})}
\begin{document}
\section*{Exchange-rate forecasting}
\textbf{Problem ID:} \texttt{F37-573}\quad \textbf{JEL:} \texttt{F37}\par
% BEGIN ECONOMICS STATEMENT
\label{problem:OP-0330}
{\sffamily\small\color{muted}Original subject group: International finance and exchange rates\par}
\label{definitions:problem:30007612}
\textbf{Audit classification:} Background; proposed extension. The classic paper reports a historical forecasting comparison, while the 2024 IMF paper presents a conditional forecasting improvement. Neither certifies the exact broad catalogue performance target as unresolved.
\subsection*{Definitions and precise research target}
Fix a prespecified set of currency pairs \(C\), forecast horizons \(H\), evaluation dates \(T\), and an information set available in real time. Let \(s_{ct}\) be the logarithm of currency pair \(c\)'s spot exchange rate. A forecasting procedure \(f\) maps dated observations to \(\widehat s^f_{c,t+h}\), with parameter estimation using only information available by \(t\). The no-drift random-walk benchmark is \(\widehat s^{RW}_{c,t+h}=s_{ct}\). For a declared loss function \(\ell\), define
\[\Delta_{ch}(f)=E[\ell(s_{c,t+h},\widehat s^f_{c,t+h})-\ell(s_{c,t+h},s_{ct})].\]
The task is to identify methods and economic conditions with negative, statistically supported loss differences on untouched observations. Specify whether the claim requires improvement for every \((c,h)\in C\times H\) or only a weighted aggregate; these are different research targets. Account for predictor revisions, model selection, overlapping forecast errors and structural breaks. Compare statistical accuracy with feasible trading gains after costs rather than treating them as identical. Derive confidence sets or predictive guarantees under stated stability assumptions and test sensitivity across regimes. A method that beats the benchmark in one sample is a resolved subcase, not a universal solution; neither the classic comparison nor the present task asserts impossibility of improved forecasting.

\textbf{Additional formal definitions and scope (editorial).} For squared-error evaluation, require \(E[(s_{c,t+h}-\widehat s^f_{c,t+h})^2]<\infty\) for each pair and horizon. An evaluation date belongs to a fixed future distribution or sampling design rather than an unspecified expectation over a historical table. A forecasting algorithm includes training windows, tuning and predictor-release calendars, and is measurable with respect to the real-time information \(\mathcal I_t\). Universal improvement means \(\Delta_{ch}<0\) for every declared \((c,h)\); weighted improvement means \(\sum_{c,h}w_{ch}\Delta_{ch}<0\), \(w_{ch}\geq0\), \(\sum w_{ch}=1\). Neither follows from a unit root test: an integrated process can have predictable stationary components, and an efficient market need not make exchange rates a no-drift random walk. Guarantees require an explicit class of stable or changing laws; unrestricted distributions can make the benchmark already optimal.

For the identification part, declare an admissible class \(\Theta\) of structural models, including preferences, technologies, shock laws, measurement/sampling rules and any equilibrium selector. If \(O\) denotes the complete observable history and \(T(\theta)\) the target defined above, the identified set is \(\mathcal I_T=\{T(\theta):\theta\in\Theta,\ \mathcal L_\theta(O)=\mathcal L_{\rm obs}(O)\}\); point identification means this set is a singleton. A \(do(a)\) comparison replaces stated policy/technology equations while fixing the declared exogenous law and re-solving the remaining equations. These definitions do not assert that an identification theorem holds without further restrictions.
\subsection*{Variants and assumption changes}
\begin{enumerate}
\item \textbf{Fixed pair and horizon (editorial).} Test a preregistered real-time forecasting rule against the no-drift benchmark for one \((c,h)\) on a fixed untouched evaluation distribution.
\item \textbf{Cross-regime predictive guarantees (editorial).} Define a class of changing laws and a weighted set of pairs/horizons; seek uniform or bounded-regret performance under those restrictions rather than unrestricted universal success.
\end{enumerate}
\subsection*{References and source audit}
\begin{enumerate}
\item Author publication record verifies Meese and Rogoff, 1983, Journal of International Economics 14:3--24; publisher indexed abstract verifies DOI/PII. Locator: Journal abstract, p.3. Support: Historical forecast comparison, not impossibility. Publisher direct fetch failed; author-hosted scan opened, without extracted text. URL: \url{https://rogoff.scholars.harvard.edu/publications/empirical-exchange-rate-models-seventies-do-they-fit-out-sample}.
\item Primary IMF WP/24/252 cover verifies Bas B. Bakker, December 2024; DOI associated with official WP. Locator: Abstract; Sections 9--10, printed pp.35--39. Support: Conditional forecasting solution. Full PDF accessible; broad exact open status not asserted. URL: \url{https://www.imf.org/-/media/files/publications/wp/2024/english/wpiea2024252-print-pdf.pdf}.
\end{enumerate}
\begin{enumerate}\raggedright
\item\hypertarget{definitions:source:30007612:1}{} Richard A. Meese; Kenneth Rogoff. 1983. \emph{Empirical Exchange Rate Models of the Seventies: Do They Fit Out of Sample?}. Journal abstract, p.3.
\url{https://rogoff.scholars.harvard.edu/publications/empirical-exchange-rate-models-seventies-do-they-fit-out-sample}.
DOI: \url{https://doi.org/10.1016/0022-1996(83)90017-X}.
\item\hypertarget{definitions:source:30007612:2}{} Bas B. Bakker. 2024. \emph{Reconciling Random Walks and Predictability: A Dual-Component Model of Exchange Rate Dynamics}. Abstract; Sections 9--10, printed pp.35--39.
\url{https://www.imf.org/-/media/files/publications/wp/2024/english/wpiea2024252-print-pdf.pdf}.
DOI: \url{https://doi.org/10.5089/9798400295034.001}.
\end{enumerate}

\subsection*{Common conventions: Source mathematical conventions}

These conventions apply to each entry unless explicitly replaced.

\textbf{Models and identification.} A primitive model \(m\) specifies all probability laws, feasible choices, information, equilibrium and measurement rules used by its target. The admissible class \(\mathcal M\) is fixed independently of outcomes. If \(P_O(m)\) is its observable probability law and \(\tau(m)\) the target, its identified set at observable law \(P\) is
\[
 \mathcal I_\tau(P)=\{\tau(m):m\in\mathcal M,\ P_O(m)=P\}.
\]
Point identification means that this set is a singleton. An empty set signals incompatibility of data and assumptions. Coordinate bounds need not describe the joint set. Bounds are sharp only if every retained target is attainable or arbitrarily approachable by compatible models. Finite bounds require explicit restrictions on unobserved quantities; unrestricted confounding, hidden wealth, extrapolation or arbitrary equilibrium selection can yield uninformative sets.

\textbf{Causal interventions.} A potential outcome \(Y_i(p)\) is the outcome under a complete policy regime \(p\), including timing, financing, information and market spillovers. The target population and units are fixed before treatment. With interference, use the complete assignment vector or a justified exposure mapping; individual no-interference notation alone cannot identify an economy-wide intervention. A difference of mean potential outcomes does not require the cross-world joint distribution. Conditioning on post-treatment migration, survival, enrollment or employment changes the estimand and needs separate justification.

\textbf{Statistical statements.} An estimator, confidence set, forecast or test specifies a sampling law, model class and loss/coverage criterion. Uniform \((1-\alpha)\)-coverage means \(\inf_{m\in\mathcal M}\Pr_m\{\tau(m)\in C_n\}\geq1-\alpha\), or an explicitly asymptotic version. The whole parameter space trivially has coverage, so informativeness requires a stated width, risk or power objective. A forecasting improvement is relative to a named benchmark, information set and evaluation population.

\textbf{Equilibrium and optimization.} An equilibrium comprises feasible plans, optimality, beliefs where needed, budget constraints and all market/resource clearing. The primitives or a stated benchmark define each correspondence \(\mathcal E(p;m)\). Nonemptiness is assumed or proved, never inferred just from compact policy sets. A selected-equilibrium target uses a declared selection rule; a robust target quantifies over all permitted equilibria. Use suprema or infima unless attainment is established. An \(\varepsilon\)-optimal policy is feasible and within \(\varepsilon\) of the finite optimum value. Report an empty feasible set separately.

\textbf{Welfare and accounting.} Normative utility, interpersonal normalization, social weights, numeraire, discounting and the use of public receipts are primitives. Behavioral utility need not coincide with the welfare criterion. Household welfare includes ownership income and rents once; adding profits again to compensating variation generally double-counts them. A monetary equivalent is defined by an explicit transfer and price convention, with monotonicity and range conditions. Average effects, mechanism contributions, transfers and welfare losses are different targets. Infinite sums require convergence and dynamic feasible plans require terminal or no-Ponzi conditions.

\textbf{Source versions.} A working-paper date, revision date and journal publication year are distinguished. A DOI for a journal version is not automatically the DOI of an earlier manuscript. Locators refer to the stated version and printed pages where specified. A metadata-only check does not verify a claim about a full-text conclusion. Every entry's source subsection narrows the provenance claim accordingly.
% END ECONOMICS STATEMENT
\end{document}
